\section{Ablation Studies}
\label{sec:ablation}

\subsection{Loss Function}
\label{ssec:ablation-loss}

Table~\ref{tab:ablation-loss} compares the dynamic Chamfer Distance loss
against displacement-based alternatives evaluated on the held-out validation
sequence (\textit{redandblack}, r2, 25 validation frames sampled during
training).

\begin{table}[t]
  \centering
  \caption{%
    Loss function ablation on \textit{redandblack} r2.
    Negative $\delta$ means degradation; ``collapse'' means the model
    predicted near-zero displacement for all points.
  }
  \label{tab:ablation-loss}
  \begin{tabular}{lcc}
    \toprule
    Loss function & Edge $\delta$ (dB) & Flat $\delta$ (dB) \\
    \midrule
    L1 displacement          & $-0.12$ & $+0.01$ (collapse) \\
    Smooth-L1 displacement   & $-0.08$ & $+0.03$ (collapse) \\
    Curvature-weighted L1    & $-0.15$ & $+0.02$ (collapse) \\
    Min-of-$k$=5 targets     & $+0.11$ & $+0.08$ \\
    Gated zero-shrinkage     & $+0.34$ & $+0.19$ \\
    Chamfer Distance (this work) & $\mathbf{+1.87}$ & $\mathbf{+1.10}$ \\
    \bottomrule
  \end{tabular}
\end{table}

All displacement-based losses produce near-zero or negative results,
consistent with the zero-inflation failure mode described in Section~\ref{sec:loss}.
The min-of-$k$ and gated zero-shrinkage approaches (Section~\ref{sec:failed-approaches})
partially address the problem but fall far short of Chamfer Distance.

\subsection{Input Features}
\label{ssec:ablation-curvature}

Table~\ref{tab:ablation-curvature} compares models trained with and without
explicit curvature features appended to the input coordinates.
The ``curv'' variant computes surface variation on the decoded cloud and
appends it as a 4th input channel.

\begin{table}[t]
  \centering
  \caption{%
    Input feature ablation. Results on 8iVFB, 4 sequences, 400 frames.
    Mean $\delta$ (dB) averaged over all sequences.
  }
  \label{tab:ablation-curvature}
  \begin{tabular}{lccccc}
    \toprule
    & \multicolumn{2}{c}{r2} & \multicolumn{2}{c}{r4} & r6 \\
    \cmidrule(lr){2-3}\cmidrule(lr){4-5}\cmidrule(lr){6-6}
    Features & Edge $\delta$ & Flat $\delta$ & Edge $\delta$ & Flat $\delta$ & Edge $\delta$ \\
    \midrule
    xyz + curvature (v10)      & $+1.62$ & $+1.15$ & $+0.46$ & $+0.64$ & $+0.17$ \\
    xyz only (v10-nocurv)      & $\mathbf{+2.06}$ & $\mathbf{+1.47}$ & $\mathbf{+0.58}$ & $\mathbf{+0.62}$ & $\mathbf{+0.27}$ \\
    \bottomrule
  \end{tabular}
\end{table}

Removing curvature improves all metrics.
The gain from removing curvature is largest at r2 ($+0.44$~dB edge), the rate
with the strongest oversmoothing.
As explained in Section~\ref{sec:formulation}, decoded-cloud curvature is
most unreliable precisely where corrections are most needed, injecting noise
into the backbone's input at the critical locations.

\subsection{Rate Conditioning Architecture}
\label{ssec:ablation-film}

Table~\ref{tab:ablation-film} compares FiLM conditioning to alternative
rate injection strategies.

\begin{table}[t]
  \centering
  \caption{%
    Rate conditioning ablation. Results at r2 and r6, averaged over 4 sequences.
  }
  \label{tab:ablation-film}
  \begin{tabular}{lcccc}
    \toprule
    & \multicolumn{2}{c}{r2 Edge $\delta$} & \multicolumn{2}{c}{r6 Edge $\delta$} \\
    Conditioning & dB & & dB & \\
    \midrule
    No conditioning (single-rate r2) & $+1.45$ & & -- &\\
    Input concatenation (rate as feature) & $+1.71$ & & $+0.12$ & \\
    FiLM on all backbone blocks & diverged & & -- &\\
    FiLM on head only (this work) & $\mathbf{+2.06}$ & & $\mathbf{+0.27}$ & \\
    \bottomrule
  \end{tabular}
\end{table}

FiLM on the head only is the best-performing approach.
Applying FiLM to all backbone blocks allows the rate signal to interfere with
the shared backbone features, causing gradient conflicts between rates during
multi-rate training and eventually diverging.
Confining rate-dependent computation to the head prevents this interference
while providing sufficient rate adaptation.

\subsection{Rate Weighting Strategy}
\label{ssec:ablation-weighting}

Table~\ref{tab:ablation-weights} compares the three rate-weighting strategies
described in Section~\ref{sec:multi-rate}.

\begin{table}[t]
  \centering
  \caption{%
    Rate weighting ablation.
    Average edge $\delta$ (dB) over 4 sequences, 400 frames.
    BD-PSNR computed over r2/r4/r6 for edge stratum.
  }
  \label{tab:ablation-weights}
  \begin{tabular}{lcccc}
    \toprule
    Weighting strategy & r2 & r4 & r6 & BD-PSNR \\
    \midrule
    Kendall uncertainty & $+0.71$ & $+0.12$ & $+0.04$ & $+0.29$ \\
    Dynamic detached    & $+1.28$ & $+0.38$ & $+0.21$ & $+0.62$ \\
    Static $[1.0, 0.3, 0.05]$ (this work) & $\mathbf{+2.06}$ & $\mathbf{+0.58}$ & $+0.27$ & $\mathbf{+0.97}$ \\
    \bottomrule
  \end{tabular}
\end{table}

Static weighting outperforms both alternatives, showing that assigning most
gradient budget to r2 is the right inductive bias.
The Kendall weighting failure is particularly informative: within 5 epochs,
the learned task weights converge to the inverted order
$[w_2, w_4, w_6] \approx [0.12, 0.31, 0.57]$, downweighting r2 because its
high loss magnitude is misinterpreted as high epistemic uncertainty.

% ============================================================
